\section{Exterior-power Husimi convex order}
\label{ext:section}

\subsection{Slater states and the measurement}

Let $E=\C^n$ with its standard Hermitian inner product, linear in the
second variable. For $0\le p\le n$, equip
\[
 \mathcal H_{p,n}=\Lambda^p E,
 \qquad D_{p,n}=\dim\mathcal H_{p,n}=\binom np,
\]
with the inner product for which increasing wedges of an orthonormal
basis are orthonormal. Thus a wedge of $p$ orthonormal vectors has norm
one. The induced unitary representation is
$R_p(U)=\Lambda^p U$, for $U\in U(n)$.

If $W\subset E$ is a $p$-dimensional subspace and $(w_1,\ldots,w_p)$
is an orthonormal basis of $W$, the vector
$\Omega_W=w_1\wedge\cdots\wedge w_p$ is a \emph{unit Slater vector}.
A different orthonormal basis changes it by a scalar of modulus one.
Consequently,
\[
 \rho_W=|\Omega_W\rangle\langle\Omega_W|
\]
is determined by $W$. We call $\rho_W$ a \emph{pure Slater density} and
write $\pi_W$ for the orthogonal projection from $E$ onto $W$; these
operators act on different spaces. The Slater-projector orbit is
$\mathcal O_{p,n}=\{\rho_W:\dim W=p\}$.

Fix a unit Slater reference $\Omega_0$. The pushforward of Haar
probability measure on $U(n)$ under
\[
 U\longmapsto
 |R_p(U)\Omega_0\rangle\langle R_p(U)\Omega_0|
\]
is denoted by $\nu_{p,n}$. Transitivity and Haar invariance make this
measure independent of the reference. Equivalently, it is the invariant
probability measure on the Grassmannian of $p$-planes, expressed through
their Slater projectors.

A \emph{density operator} on $\mathcal H_{p,n}$ is an operator
$\rho\ge0$ with $\Tr\rho=1$. Its Husimi function is
\begin{equation}
 Q_\rho(P)=\Tr(\rho P),\qquad P\in\mathcal O_{p,n}.
 \label{ext:husimi}
\end{equation}
We also write
$Q_\rho(U)=\langle R_p(U)\Omega_0,\rho R_p(U)\Omega_0\rangle$
for the same measurement pulled back to the group. Since
$0\le\rho\le\Id$, its values lie in $[0,1]$.

The normalization is
\begin{equation}
 \int_{\mathcal O_{p,n}}P\,d\nu_{p,n}(P)
 =\frac{\Id_{\mathcal H_{p,n}}}{D_{p,n}},
 \qquad
 \int Q_\rho\,d\nu_{p,n}=\frac1{D_{p,n}}.
 \label{ext:resolution}
\end{equation}
Indeed, the averaged projector commutes with every $R_p(U)$.
Commutation with diagonal one-particle unitaries makes it diagonal in
the occupation basis, because distinct occupied subsets have distinct
characters. Permutation unitaries make all its diagonal entries equal.
Its trace is one, which proves the first identity; pairing with $\rho$
proves the second. For $p=0$ and $p=n$, the sector is one-dimensional,
its unique density is Slater, and $Q_\rho\equiv1$. This also includes
$n=p=0$.

\begin{theorem}[Exterior-power Husimi convex order]
\label{ext:convex-order}
For every $0\le p\le n$, every density operator $\rho$ on
$\mathcal H_{p,n}$, every continuous convex function
$\Phi:[0,1]\to\R$, and every pure Slater density $\rho_W$,
\begin{equation}
 \int\Phi(Q_\rho)\,d\nu_{p,n}
 \le
 \int\Phi(Q_{\rho_W})\,d\nu_{p,n}.
 \label{ext:order}
\end{equation}
The right-hand side is independent of $W$. If $\Phi$ is strictly
convex, equality holds if and only if $\rho$ is a pure Slater density.
\end{theorem}

All the compared distributions have the mean in
\eqref{ext:resolution}. The theorem says that a Slater input maximizes
every continuous convex observation of the measurement value. The
equality statement concerns the whole density operator, including its
off-diagonal entries.

\subsection{A shared convex mixture}

We first record the probabilistic step used here and in the spinor
argument. Its coefficients are fixed before the measurement variable
is sampled.

\begin{lemma}[Shared-mixture comparison]
\label{ext:common-mixture}
Let $(X,\mu)$ be a probability space. Suppose measurable functions
$a,a_1,\ldots,a_L:X\to[0,1]$ satisfy
\[
 a(x)=\sum_{j=1}^L t_j a_j(x),\qquad
 t_j\ge0,\qquad\sum_{j=1}^L t_j=1,
\]
and all $a_j$ have the same distribution $\eta$. For every continuous
convex function $H:[0,1]\to\R$,
\begin{equation}
 \int_X H(a(x))\,d\mu(x)
 \le \int_{[0,1]}H(t)\,d\eta(t).
 \label{ext:shared-jensen}
\end{equation}
If $H$ is strictly convex, equality implies that all the $a_j$ with
$t_j>0$ agree almost everywhere.
\end{lemma}
\begin{proof}
Pointwise Jensen gives $H(a)\le\sum_jt_jH(a_j)$. Integration and
the common distribution give \eqref{ext:shared-jensen}. All functions
being integrated are bounded. If the integrals are equal, the
nonnegative Jensen difference vanishes almost everywhere. The equality
criterion for strict convexity gives the assertion.
\end{proof}

If $R$ is a $[0,1]$-valued random variable, define
\begin{equation}
 H_\Phi(t)=\E\Phi(tR),\qquad 0\le t\le1.
 \label{ext:scale-function}
\end{equation}
Continuity follows by bounded convergence, and convexity follows
pointwise. If $\Phi$ is strictly convex and $\Pr(R>0)>0$, then
$H_\Phi$ is strictly convex: for distinct $s,t$ the strict scalar
inequality holds on the event $R>0$.

\subsection{Contraction and the one-particle density}

For $u\in E$, let $c(u)^*$ denote creation by $u$, namely
$\xi\mapsto u\wedge\xi$, and let $c(u)$ be its adjoint. Thus $c(u)$
is conjugate-linear in $u$. These operators act on the full graded
exterior space $\Lambda^\bullet E$ and change degree by one; contraction
vanishes on degree zero. On a decomposable vector,
\begin{equation}
 c(u)(v_1\wedge\cdots\wedge v_p)
 =\sum_{j=1}^p(-1)^{j-1}\langle u,v_j\rangle
 v_1\wedge\cdots\wedge\widehat v_j\wedge\cdots\wedge v_p.
 \label{ext:contraction}
\end{equation}
The hat indicates the omitted factor. The canonical anticommutation
relations are
\[
 c(u)c(v)+c(v)c(u)=0,
 \qquad
 c(u)c(v)^*+c(v)^*c(u)=\langle u,v\rangle\Id.
\]
For unit $u$, the operator $N_u=c(u)^*c(u)$ is the projection onto
the states in which $u$ is occupied. If $(e_i)_{i=1}^n$ is an
orthonormal basis, then on $\mathcal H_{p,n}$,
\begin{equation}
 \sum_{i=1}^n N_{e_i}=p\Id.
 \label{ext:number}
\end{equation}
Both assertions can be read directly in the occupation basis.

For a density $\rho$ on $\mathcal H_{p,n}$, define its one-particle
density $\gamma_\rho$ by
\begin{equation}
 a_\rho(u):=\Tr\bigl(c(u)\rho c(u)^*\bigr)
 =\Tr(\rho N_u)=\langle u,\gamma_\rho u\rangle.
 \label{ext:one-body}
\end{equation}
Polarization gives a unique Hermitian operator on $E$. The projection
property of $N_u$ and \eqref{ext:number} imply
\begin{equation}
 0\le\gamma_\rho\le\Id_E,
 \qquad \Tr\gamma_\rho=p.
 \label{ext:pauli-bounds}
\end{equation}
For a Slater input, $\gamma_{\rho_W}=\pi_W$.

\begin{lemma}[One-particle mixture from the original density]
\label{ext:occupation-mixture}
Choose an orthonormal eigenbasis $(f_i)_{i=1}^n$ of $\gamma_\rho$.
For every $p$-element subset $S\subset\{1,\ldots,n\}$, let
$\Omega_S$ be its increasing occupation wedge,
$W_S=\operatorname{span}\{f_i:i\in S\}$, and
$t_S=\langle\Omega_S,\rho\Omega_S\rangle$. Then
\begin{equation}
 t_S\ge0,\qquad \sum_{|S|=p}t_S=1,
 \qquad
 \gamma_\rho=\sum_{|S|=p}t_S\pi_{W_S}.
 \label{ext:diagonal-mixture}
\end{equation}
If $n\ge1$, then for a uniformly distributed unit vector $u\in E$,
any continuous convex $H:[0,1]\to\R$, and any rank-$p$ projection $\pi_W$,
\begin{equation}
 \E_u H(\langle u,\gamma_\rho u\rangle)
 \le \E_u H(\langle u,\pi_Wu\rangle).
 \label{ext:occupation-order}
\end{equation}
If $H$ is strictly convex, equality forces exactly one $t_S$ to be
positive, and consequently $\rho=|\Omega_S\rangle\langle\Omega_S|$.
\end{lemma}
\begin{proof}
The $t_S$ are diagonal entries of the same positive trace-one operator
in a complete orthonormal basis, so they are probabilities. For each
$i$, the occupation operator $N_{f_i}$ is diagonal in this basis and
has eigenvalue $1$ precisely when $i\in S$. Hence
\[
 \langle f_i,\gamma_\rho f_i\rangle
 =\Tr(\rho N_{f_i})=\sum_{S\ni i}t_S.
\]
The off-diagonal entries of $\gamma_\rho$ vanish in the chosen
eigenbasis. This proves \eqref{ext:diagonal-mixture}.
In particular, for every unit $u$ at once,
\[
 a_\rho(u)=\sum_{|S|=p}t_S\langle u,\pi_{W_S}u\rangle.
\]
The projections on the right have the same rank, so their quadratic
readouts have the same spherical distribution. Lemma~\ref{ext:common-mixture}
gives \eqref{ext:occupation-order}.

For $S\ne T$, the Hermitian operator $\pi_{W_S}-\pi_{W_T}$ is
nonzero. Its quadratic form is nonzero somewhere on the unit sphere,
and therefore on a nonempty open set of positive spherical measure.
The strict equality criterion excludes two such positive-weight
indices. The remaining diagonal entry of $\rho$ is one. Positivity
gives
\[
 |\langle\Omega_S,\rho\Omega_T\rangle|^2
 \le t_S t_T
\]
for all occupied subsets, so every other matrix entry is zero.
Thus the original density is the stated pure Slater projector.
\end{proof}

\begin{remark}[The Pauli convex set]
\label{ext:pauli-convex-set}
The set of Hermitian matrices satisfying \eqref{ext:pauli-bounds} is
the convex hull of rank-$p$ orthogonal projections. To see this,
diagonalize such a matrix. The vertices of
$\{x\in[0,1]^n:\sum_i x_i=p\}$ are the zero-one vectors with $p$
ones: a nonintegral point has at least two nonintegral coordinates,
which admit opposite small perturbations preserving the sum.
Rotating a vertex decomposition back proves the assertion.
For the density in the theorem, \eqref{ext:diagonal-mixture} supplies
this decomposition directly from its occupation probabilities and
also yields the full-density equality criterion.
\end{remark}

\subsection{The conditional measurement}

For a unit vector $u$, put
$A_\rho(u)=c(u)\rho c(u)^*$. It is positive with trace
$a_\rho(u)$. The orthogonal splitting
\[
 \Lambda^{p-1}E
 =\Lambda^{p-1}(u^\perp)
 \mathbin{\oplus}
 u\wedge\Lambda^{p-2}(u^\perp)
\]
shows that $\ker c(u)$ in degree $p-1$ is
$\Lambda^{p-1}(u^\perp)$; the second summand is absent when $p=1$.
Since $c(u)^2=0$, the range of $A_\rho(u)$ lies in this kernel.
If $a_\rho(u)>0$, then
\begin{equation}
 \rho_u=\frac{A_\rho(u)}{a_\rho(u)}
 \label{ext:conditional-density}
\end{equation}
is a density on $\Lambda^{p-1}(u^\perp)$.
For a unit Slater vector $\xi$ in that space, adjointness gives
\begin{equation}
 \langle u\wedge\xi,\rho(u\wedge\xi)\rangle
 =\langle\xi,A_\rho(u)\xi\rangle
 =a_\rho(u)\langle\xi,\rho_u\xi\rangle.
 \label{ext:conditioning}
\end{equation}
If $a_\rho(u)=0$, positivity and zero trace give $A_\rho(u)=0$,
so the first expression in \eqref{ext:conditioning} is zero as well.

The relevant Haar decomposition can be written without choosing a
global frame for the bundle $u^\perp$. For $p\ge1$, take
$\Omega_0=e_1\wedge\cdots\wedge e_p$ and embed $U(n-1)$ into
$U(n)$ by $\iota(V)=1\oplus V$. Right invariance gives
\begin{equation}
 \int_{U(n)}\Phi(Q_\rho(U))\,dU
 =\int_{U(n)}\int_{U(n-1)}
 \Phi\bigl(Q_\rho(U\iota(V))\bigr)\,dV\,dU.
 \label{ext:double-haar}
\end{equation}
For fixed $U$, the first vector is $u=Ue_1$, and the remaining wedge
is a Haar Slater vector in $\Lambda^{p-1}(u^\perp)$. The outer
variable $u$ is uniform on the complex unit sphere. Thus
\eqref{ext:conditioning} applies to the same measurement appearing
in \eqref{ext:double-haar}. In particular, the outer measure remains
Haar probability measure; $a_\rho(u)$ enters the argument of $\Phi$.

\begin{lemma}[Slater closure under contraction]
\label{ext:coherent-closure}
If $a_{\rho_W}(u)>0$, the conditional density $(\rho_W)_u$ is a
pure Slater density with occupied space $W\cap u^\perp$.
\end{lemma}
\begin{proof}
Here $a_{\rho_W}(u)=\|\pi_Wu\|^2$. Choose an orthonormal basis of
$W$ whose first vector is $w_1=\pi_Wu/\|\pi_Wu\|$.
Its other vectors $w_2,\ldots,w_p$ lie in $u^\perp$.
Up to the phase of $\Omega_W$, formula~\eqref{ext:contraction} gives
\[
 c(u)\Omega_W
 =\|\pi_Wu\|\,w_2\wedge\cdots\wedge w_p.
\]
Normalizing its rank-one projector proves the assertion, including
the vacuum conditional state when $p=1$.
\end{proof}

\subsection{Induction and equality}

\begin{proof}[Proof of Theorem~\ref{ext:convex-order}]
We induct on $p$, simultaneously for all $n\ge p$. The case $p=0$
has $Q\equiv1$. Suppose $p\ge1$ and the assertion is known in
degree $p-1$. Let $R$ have the Husimi distribution of a Slater input
in $\mathcal H_{p-1,n-1}$ and set $H_\Phi(t)=\E\Phi(tR)$.
This function is continuous and convex. If $\Phi$ is strictly
convex, then $H_\Phi$ is strictly convex, since
\[
 \E R=\binom{n-1}{p-1}^{-1}>0
\]
by \eqref{ext:resolution}.

For every positive-weight conditional density \eqref{ext:conditional-density},
apply the induction hypothesis with the test function
$x\mapsto\Phi(a_\rho(u)x)$. Zero-weight branches give $\Phi(0)$
on both sides. Integrating \eqref{ext:double-haar}, and then applying
Lemma~\ref{ext:occupation-mixture}, yields
\begin{align}
 \int\Phi(Q_\rho)\,d\nu_{p,n}
 &\le\E_u H_\Phi(\langle u,\gamma_\rho u\rangle)\notag\\
 &\le\E_u H_\Phi(\langle u,\pi_Wu\rangle)
 =\int\Phi(Q_{\rho_W})\,d\nu_{p,n}.
 \label{ext:chain}
\end{align}
The final equality follows from Lemma~\ref{ext:coherent-closure}:
each nonzero conditional Slater input has the same lower-dimensional
Husimi distribution. The original double-Haar integrands are continuous
on compact groups. The normalized conditional matrices are measurable
where their continuous weights are positive, and the zero branch was
handled separately. Boundedness therefore justifies these integrations
and their order changes.

If $\Phi$ is strictly convex and the first and last quantities in
\eqref{ext:chain} are equal, the middle Jensen comparison is an
equality. Lemma~\ref{ext:occupation-mixture}, applied to the strictly
convex $H_\Phi$, forces $\rho$ to be a pure Slater density.
Conversely, all Slater inputs have the same Husimi law by Haar
invariance and hence attain equality. The filled sector $p=n$ is
one-dimensional and agrees with this induction.
\end{proof}

\begin{example}[One particle]
\label{ext:one-particle-example}
Every unit vector in $E$ is a one-particle Slater vector. In two
modes, a density with eigenvalues $t$ and $1-t$ has measurement
\[
 Q_\rho(u)=t|u_1|^2+(1-t)|u_2|^2.
\]
At $t=1/2$ this is the constant $1/2$; at $t=1$ it has the
distribution of $|u_1|^2$. Both have mean $1/2$. Theorem~\ref{ext:convex-order}
compares all continuous convex observations between these cases and
every intermediate density. Its strict equality criterion singles
out the rank-one inputs.
\end{example}

\begin{example}[Two-particle coherence]
\label{ext:two-particle-example}
In $\Lambda^2\C^4$, write $e_{ij}=e_i\wedge e_j$ and consider
\[
 \psi=\frac{e_{12}+e_{34}}{\sqrt2},\qquad
 \rho_{\mathrm{mix}}
 =\tfrac12|e_{12}\rangle\langle e_{12}|
  +\tfrac12|e_{34}\rangle\langle e_{34}|.
\]
Both $|\psi\rangle\langle\psi|$ and $\rho_{\mathrm{mix}}$ have
one-particle density $\tfrac12\Id_E$. In the standard eigenbasis,
the shared occupation mixture places weight $1/2$ on each of
$\{1,2\}$ and $\{3,4\}$. Their full Husimi functions differ:
at the plane spanned by $(e_1+e_3)/\sqrt2$ and
$(e_2+e_4)/\sqrt2$, their values are respectively $1/2$ and $1/4$.
Thus the one-particle mixture retains exactly the information needed
for the common conditional-weight comparison, while the conditional
density in \eqref{ext:conditional-density} still carries the original
coherence. Both inputs give strict inequality in
\eqref{ext:order} for every strictly convex test.
\end{example}
